\documentclass{beamer}

\usetheme{Madrid}
\usecolortheme{default}
\usepackage{hyperref}

\title[Missing Data]{Missing Data and Multiple Imputation}
\subtitle{A Practical Introduction with \texttt{mice}}

\author{Monnie McGee}

\institute{
REU Site: Advancing Data Science in the Public Interest\\[0.5em]
Based on examples from\\
\emph{Flexible Imputation of Missing Data} (2nd ed.)\\
Stef van Buuren (2018)
}

\date{}
\begin{document}

\begin{frame}
  \titlepage
\end{frame}

\begin{frame}{Sources and Acknowledgments}

This lesson is adapted from:

\vspace{0.5em}

\textbf{Stef van Buuren (2018)}

\emph{Flexible Imputation of Missing Data}, Second Edition.

\vspace{0.5em}

Online version:

\url{https://stefvanbuuren.name/fimd/}

\vspace{0.5em}

Source code and examples adapted from:

\url{https://github.com/stefvanbuuren/fimdbook}

\vspace{1em}

Examples have been modernized using:

\begin{itemize}
\item \texttt{mice}
\item \texttt{naniar}
\item \texttt{ggplot2}
\item \texttt{VIM}
\end{itemize}

\end{frame}



\begin{frame}{Learning Objectives}
By the end of this session, you should be able to:
\begin{itemize}
  \item Explain why missing data are a statistical problem
  \item Distinguish MCAR, MAR, and MNAR
  \item Recognize limitations of ad hoc methods
  \item Understand the basic idea of multiple imputation
  \item Use the \texttt{mice} package for a simple analysis
\end{itemize}
\end{frame}

\begin{frame}{Why Should We Care About Missing Data?}
Missing data occur everywhere:
\begin{itemize}
  \item Surveys
  \item Medical studies
  \item Environmental monitoring
  \item Business data
  \item Social science research
\end{itemize}

\vspace{0.2cm}

Missing data can lead to:
\begin{itemize}
  \item Smaller sample sizes
  \item Biased estimates
  \item Reduced power
  \item Incorrect conclusions
\end{itemize}
\end{frame}

\begin{frame}{A Simple Example}
Suppose we observe:

\[
1,\; 2,\; 4
\]

\[
\bar{x} = 2.33
\]

\vspace{0.4cm}

Now suppose we observe:

\[
1,\; 2,\; \texttt{NA}
\]

\vspace{0.3cm}

\begin{center}
\Large What should the mean be?
\end{center}
\end{frame}

\begin{frame}[plain]
\centering
\vfill
{\Huge Quarto Demonstration}

\vspace{0.5cm}

{\Large Mean calculations with missing values}

\vfill
\end{frame}

\begin{frame}{Complete-Case Analysis}
Also called \textbf{listwise deletion}.

\vspace{0.3cm}

\textbf{Idea:} remove any row with missing values.

\vspace{0.3cm}

Advantages:
\begin{itemize}
  \item Easy
  \item Built into most software
\end{itemize}

Disadvantages:
\begin{itemize}
  \item Data loss
  \item Potential bias
  \item Changing sample sizes across models
\end{itemize}
\end{frame}

\begin{frame}[fragile]{Complete-Case Analysis in R}
\begin{verbatim}
lm(Ozone ~ Wind, data = airquality)

lm(Ozone ~ Wind + Solar.R, data = airquality)
\end{verbatim}

\vspace{0.4cm}

Adding a variable can change which rows are included.
\end{frame}

\begin{frame}[plain]
\centering
\vfill
{\Huge Quarto Demonstration}

\vspace{0.5cm}

{\Large Listwise deletion in \texttt{airquality}}

\vfill
\end{frame}

\begin{frame}{Types of Missing Data}
Three missingness mechanisms:

\begin{enumerate}
  \item MCAR
  \item MAR
  \item MNAR
\end{enumerate}

\vspace{0.4cm}

These mechanisms help determine whether a missing-data method is reasonable.
\end{frame}

\begin{frame}{MCAR}
\textbf{Missing Completely At Random}

\vspace{0.3cm}

Missingness is unrelated to:
\begin{itemize}
  \item observed values
  \item unobserved values
\end{itemize}

\vspace{0.3cm}

Example:

\begin{quote}
A sensor randomly fails.
\end{quote}

Observed cases are still representative.
\end{frame}

\begin{frame}{MAR}
\textbf{Missing At Random}

\vspace{0.3cm}

Missingness depends on observed variables.

\vspace{0.3cm}

Example:

\begin{quote}
Older participants are less likely to answer a question, and age is observed.
\end{quote}

\vspace{0.3cm}

Multiple imputation often works well under MAR.
\end{frame}

\begin{frame}{MNAR}
\textbf{Missing Not At Random}

\vspace{0.3cm}

Missingness depends on the missing value itself.

\vspace{0.3cm}

Example:

\begin{quote}
People with very high income are less likely to report income.
\end{quote}

\vspace{0.3cm}

This is the hardest case.
\end{frame}

\begin{frame}{Think-Pair-Share}
Classify each as MCAR, MAR, or MNAR:

\begin{enumerate}
  \item A survey website crashes randomly.
  \item Older respondents skip a question.
  \item Students with low grades avoid reporting GPA.
\end{enumerate}
\end{frame}

\begin{frame}{Why MCAR Matters}
Pairwise deletion combines statistics from different subsets.

\vspace{0.3cm}

For example:
\begin{itemize}
  \item Mean of \texttt{Ozone} may use one set of rows
  \item Mean of \texttt{Solar.R} may use another set
  \item Covariance may use yet another set
\end{itemize}

\vspace{0.3cm}
This only works well if those subsets are representative.

\vspace{0.3cm}
That is the MCAR assumption.
\end{frame}

\begin{frame}[plain]
\centering
\vfill
{\Huge Quarto Demonstration}

\vspace{0.5cm}

{\Large Pairwise deletion and sample sizes}

\vfill
\end{frame}

\begin{frame}{Common Ad Hoc Methods}
Common quick fixes:

\begin{itemize}
  \item Listwise deletion
  \item Pairwise deletion
  \item Mean imputation
  \item Regression imputation
  \item Last observation carried forward
\end{itemize}

\vspace{0.3cm}
They are simple, but often misleading.
\end{frame}

\begin{frame}{Mean Imputation}
Replace missing values with the mean.

\[
x_{\text{missing}} = \bar{x}_{\text{observed}}
\]

\vspace{0.3cm}

Advantages:
\begin{itemize}
\item Easy
\item Fast
\item Available in almost every software package
\end{itemize}

\vspace{0.3cm}

Problems:
\begin{itemize}
\item Reduces variance
\item Distorts distributions
\item Underestimates uncertainty
\item Weakens relationships among variables
\end{itemize}
\end{frame}

\begin{frame}[plain]
\centering
\vfill
{\Huge Quarto Demonstration}

\vspace{0.5cm}

{\Large Mean Imputation}

\vspace{0.5cm}

Histogram and scatterplot examples

\vfill
\end{frame}

\begin{frame}{Regression Imputation}

Fit a regression model:

\[
Y = \beta_0 + \beta_1 X
\]

Use the fitted model to predict missing values.

\vspace{0.3cm}

Advantages:

\begin{itemize}
\item Uses relationships among variables
\item Better than mean imputation
\end{itemize}

\vspace{0.3cm}

Problem:

Predicted values fall directly on the regression line.

Variability is lost.
\end{frame}

\begin{frame}[plain]
\centering
\vfill
{\Huge Quarto Demonstration}

\vspace{0.5cm}

{\Large Regression Imputation}

\vfill
\end{frame}

\begin{frame}{Last Observation Carried Forward (LOCF)}

Idea:

Replace a missing value by the most recent observed value.

\vspace{0.3cm}

Common in longitudinal studies.

\vspace{0.3cm}

Potential problems:

\begin{itemize}
\item Assumes no change occurred
\item Can create artificial trends
\item Often underestimates uncertainty
\end{itemize}
\end{frame}

\begin{frame}[plain]
\centering
\vfill
{\Huge Quarto Demonstration}

\vspace{0.5cm}

{\Large LOCF Example}

\vfill
\end{frame}

\begin{frame}{Why Ad Hoc Methods Fail}

Common problems:

\begin{itemize}
\item Loss of information
\item Biased estimates
\item Distorted distributions
\item Incorrect standard errors
\item Underestimated uncertainty
\end{itemize}

\vspace{0.5cm}

We need a method that uses available information while acknowledging uncertainty.
\end{frame}

\begin{frame}{Multiple Imputation}

Core idea:

\vspace{0.3cm}

Instead of filling in one value,

\vspace{0.2cm}

fill in several plausible values.

\vspace{0.5cm}

Create:

\begin{enumerate}
\item Multiple completed datasets
\item Separate analyses
\item Combined results
\end{enumerate}
\end{frame}

\begin{frame}[fragile]{The Multiple Imputation Workflow}

\begin{center}
\includegraphics[scale=.15]{ch01-miflow-1.png}
\end{center}

\end{frame}

\begin{frame}{What Produces the Imputed Values?}

Multiple imputation is a framework.

\vspace{0.3cm}

An imputation method generates the replacement values.

\vspace{0.3cm}

Examples available in \texttt{mice}:

\begin{itemize}
\item Predictive mean matching (PMM)
\item Bayesian regression
\item Classification and Regression Trees (CART)
\item Random Forest
\item Logistic regression
\end{itemize}
\end{frame}

\begin{frame}{Predictive Mean Matching (PMM)}

PMM is one of the most popular methods for continuous variables.

\vspace{0.3cm}

Basic idea:

\begin{enumerate}
\item Predict the missing value.
\item Find observed cases with similar predictions.
\item Randomly select one observed value.
\end{enumerate}

\vspace{0.3cm}

The imputed value is always an observed value from the data.
\end{frame}

\begin{frame}{PMM Example}

Suppose a student's GPA is missing.

\vspace{0.3cm}

Model prediction:

\[
3.72
\]

\vspace{0.3cm}

Closest observed GPAs:

\[
3.70,\quad 3.80
\]

\vspace{0.3cm}

PMM randomly selects one of those observed values.

\vspace{0.3cm}

Possible imputation:

\[
3.70
\]
or
\[
3.80
\]
\end{frame}

\begin{frame}{Why PMM Is Popular}

Advantages:

\begin{itemize}
\item Preserves realistic values
\item Preserves distributions
\item Avoids impossible values
\item Often performs well in practice
\end{itemize}

\vspace{0.5cm}

For many continuous variables, PMM is a sensible default.
\end{frame}

\begin{frame}{Common \texttt{mice} Methods}

\begin{tabular}{ll}
\textbf{Method} & \textbf{Purpose} \\
\hline
\texttt{pmm} & Predictive Mean Matching \\
\texttt{norm} & Bayesian regression \\
\texttt{norm.nob} & Stochastic regression \\
\texttt{norm.predict} & Deterministic regression \\
\texttt{cart} & Decision trees \\
\texttt{rf} & Random forest \\
\texttt{logreg} & Binary variables \\
\texttt{polyreg} & Multicategory variables \\
\texttt{polr} & Ordered categories \\
\end{tabular}
\end{frame}

\begin{frame}[plain]
\centering
\vfill
{\Huge Quarto Demonstration}

\vspace{0.5cm}

{\Large Using \texttt{mice}}

\vspace{0.5cm}

\texttt{md.pattern()}, \texttt{mice()},
\texttt{densityplot()}, \texttt{stripplot()}

\vfill
\end{frame}

\begin{frame}[fragile]{Typical \texttt{mice} Workflow}

\begin{verbatim}
imp <- mice(
data,
method = "pmm",
m = 20,
seed = 2026
)

fit <- with(
imp,
lm(y ~ x1 + x2)
)

pool(fit)
\end{verbatim}
\end{frame}

\begin{frame}{What Should You Report?}

When using multiple imputation:

\begin{itemize}
\item Number of imputations
\item Imputation method
\item Variables used in imputation
\item Analysis model
\end{itemize}

\vspace{0.5cm}

Transparency is important.
\end{frame}

\begin{frame}{Key Takeaways}

\begin{enumerate}
\item Missing data are a statistical problem.
\item Missingness mechanisms matter.
\item Ad hoc methods often distort results.
\item Multiple imputation accounts for uncertainty.
\item \texttt{mice} provides practical tools for implementation.
\end{enumerate}
\end{frame}



\begin{frame}{Further Reading}

\textbf{Primary Reference}

\vspace{0.5em}

Stef van Buuren (2018)

\emph{Flexible Imputation of Missing Data}, 2nd Edition

\vspace{0.5em}

Book website:

\url{https://stefvanbuuren.name/fimd/}

\vspace{0.5em}

\textbf{R Packages}

\begin{itemize}
\item \texttt{mice}: \url{https://amices.org/mice/}
\item \texttt{naniar}: \url{https://naniar.njtierney.com/}
\item \texttt{VIM}: \url{https://cran.r-project.org/package=VIM}
\end{itemize}

\end{frame}

\begin{frame}
\centering
\Huge Questions?
\end{frame}

\end{document}

%One suggestion: when you discuss PMM, don't present it as "the best" method. Tell students:

%> PMM is often a good default for continuous variables because it tends to preserve realistic values and is relatively robust.

%That statement is accurate, easy to remember, and avoids getting drawn into debates about Bayesian regression, random forests, or CART during an introductory lesson.
